\solchapter{bio}

\sol{bio}{1} $\kB T = 4.3\times10^{-21}\un{J} = 4.3\un{pN\,nm} = 27\un{meV}$. Kinesin does at most $8\un{nm}\times6\un{pN} = 48\un{pN\,nm}$ of work per step, from $20\kB T = 86\un{pN\,nm}$: an efficiency of about $56\%$, better than a car engine.

\sol{bio}{2} With $t\sim L^2/D$: $0.1\un{s}$, $40\un{s}$, and $10^{11}\un{s}$, or $3000$ years. Kinesin covers $1\un{m}$ in $10^6\un{s}$, about $12$ days (the slowest axonal transport takes months). Active transport wins when $L/v<L^2/D$, for $L>D/v = 10\un{\mu m}$: about the size of a cell.

\sol{bio}{3} $6\times10^9\times0.34\un{nm} = 2\un{m}$, about $3\times10^5$ times the nucleus's diameter. As a random coil of $N = 2\un{m}/100\un{nm} = 2\times10^7$ segments, its size would be $\sqrt Nb\approx0.45\un{mm}$, still $75$ times too large. DNA is actively packaged: wound around histone proteins, then coiled and looped into chromatin, in a way that still lets the cell read any gene on demand.

\sol{bio}{4} $V = (\kB T/e)\ln(c_{\text{out}}/c_{\text{in}})$, with $\kB T/e = 26.7\un{mV}$. Potassium: $26.7\ln(5/140) = -89\un{mV}$. Sodium: $26.7\ln(145/12) = +67\un{mV}$. At rest the membrane is mostly permeable to potassium, so its potential sits near $-70\un{mV}$, close to the potassium value. When sodium channels open, sodium becomes the dominant permeability, and it rushes in down both its concentration gradient and the electric field, driving the potential toward $+67\un{mV}$: the spike.

\sol{bio}{5} $S = -3\kB R^2/2Nb^2 + \text{const}$, so $F = -T\,\partial S/\partial R = 3\kB TR/Nb^2$: the force needed to hold the ends a distance $R$ apart. The polymer behaves as a Hookean spring of stiffness $3\kB T/Nb^2$, which pulls the ends together. The force is entirely entropic: stretching reduces the number of available configurations. The stiffness is proportional to $T$ because the free energy cost of losing entropy is $T\Delta S$. A heated rubber band contracts, unlike almost any other material --- try it with a hair dryer and a weight.

\sol{bio}{6} The drag coefficient is $6\pi\eta r = 6\pi(10^{-3})(0.5\times10^{-6}) = 9.4\times10^{-9}\un{kg/s}$, so the stopping time is $m/6\pi\eta r = 10^{-7}\un{s}$, and the coasting distance $v\tau = 3\times10^{-12}\un{m}$: a thirtieth of the diameter of a hydrogen atom. A bacterium lives at a Reynolds number of $10^{-5}$, where inertia is irrelevant: it stops dead the instant it stops swimming.

\sol{bio}{7} $3^{100}\approx5\times10^{47}$ conformations, which at $10^{12}$ per second takes $5\times10^{35}\un{s}$, or $10^{28}$ years --- $10^{18}$ times the age of the universe. A protein that folds in $1\un{s}$ can sample only about $10^{12}$ conformations, so the effective search space must be reduced by a factor of about $10^{35}$ (or $10^{41}$ for microsecond folders). Folding is not a random search: the energy landscape is a funnel, and the protein slides downhill, forming local structures that guide the rest.

\sol{bio}{8} Metabolic rate per kilogram scales as $M^{-1/4}$, so the mouse's is $(5000/0.02)^{1/4} = 22$ times the elephant's. A mouse at $600$ beats per minute uses up $10^9$ heartbeats in $3.2$ years (real mice live $2$--$3$ years). If heart rate also scales as $M^{-1/4}$, the elephant's is $600/22 = 27$ beats per minute (measured: about $30$), and it lives about $70$ years (real elephants: $60$--$70$).

\sol{bio}{9} A squid's giant axon conducts at about $25\un{m/s}$. Myelinated axons conduct at roughly $6\un{m/s}$ per micrometre of diameter, so a myelinated axon about $4\un{\mu m}$ across matches it: some $100$ times thinner, with $10^4$ times less cross-section. A vertebrate nerve can pack thousands of fast fibres into the space of one squid axon, which is what makes complex nervous systems --- and big brains --- affordable.

\sol{bio}{10} A molecule near the receptor wanders away in a time $\sim a^2/D$, and in that time the receptor holds about $ca^3$ molecules. In a time $T$ it gets $TD/a^2$ independent looks, so it counts about $ca^3\cdot TD/a^2 = DacT$ molecules altogether --- the same as the diffusive arrival rate $\sim Dac$ times $T$. Poisson statistics give an error $\sqrt N$ in a count $N$, so $\delta c/c\sim1/\sqrt{DacT}$. For a small molecule, $D\approx10^{-9}\un{m^2/s}$, and $c = 1\un{\mu M} = 6\times10^{20}\un{m^{-3}}$. With $a = 1\un{\mu m}$ and $T = 1\un{s}$, $DacT\approx6\times10^5$ and $\delta c/c\approx0.1\%$. Measuring for longer helps, but a bacterium swimming in a straight line is knocked off course by rotational diffusion within about $10\un{s}$, so it cannot average much longer; and it must swim far enough in that time to outrun diffusion, or it could not tell where the concentration was higher.

\sol{bio}{11} Suppose a forward step requires the gas to supply energy $\epsilon$ (to lift the pawl over a tooth) plus $L\theta$ (to lift the weight), with probability $\propto e^{-(\epsilon+L\theta)/\kB T_1}$ per attempt. A backward step happens when the pawl, at temperature $T_2$, is thermally lifted by $\epsilon$ (probability $\propto e^{-\epsilon/\kB T_2}$), letting the weight turn the ratchet back. With no load and $T_1 = T_2$, the rates are equal: no net rotation, as the second law requires. In general, the ratchet turns forward when $(\epsilon+L\theta)/T_1<\epsilon/T_2$. Each forward step takes heat $Q_1 = \epsilon+L\theta$ from the gas, does work $W = L\theta$, and delivers $Q_2 = \epsilon$ to the pawl (when it snaps down). Its efficiency, $W/Q_1 = 1 - \epsilon/(\epsilon+L\theta)$, is less than $1 - T_2/T_1$ exactly when it runs forward: never better than Carnot, and equal to Carnot when the device is balanced and runs infinitely slowly. Kinesin does not run on heat at all. It runs on chemical free energy: the cell keeps ATP far from equilibrium with ADP and phosphate, and each hydrolysis, by changing the motor's shape, switches the potential landscape the motor head sits in. Thermal motion moves the head; the irreversible chemical cycle rectifies it, paying about $20\kB T$ a step.

\sol{bio}{12} If the right and wrong substrates bind at the same rate but the wrong one falls off faster by $e^{\Delta G/\kB T}$, equilibrium occupancy discriminates by $e^{-\Delta G/\kB T}$. After the activation step, the complex either makes product at rate $W$ or releases the substrate, at a rate $\ell'$ that again is $e^{\Delta G/\kB T}$ times larger for the wrong one. If $W\ll\ell'$, the fraction making product is $W/\ell'$, a second factor of $e^{-\Delta G/\kB T}$: the error rate becomes $e^{-2\Delta G/\kB T}$. A substrate that falls off the activated complex cannot climb directly back into it: it must start over with binding. This one-way flow is the point. If the activation step were reversible and at equilibrium, the activated complex would be just one more bound state, whose population is fixed by its free energy alone --- and equilibrium can never discriminate by more than $e^{-\Delta G/\kB T}$, however many states there are. Driving the cycle by GTP or ATP hydrolysis breaks detailed balance. The price is energy, and some correct substrates discarded: ribosomes and DNA polymerases trade both for accuracy.
